% Options for packages loaded elsewhere
\PassOptionsToPackage{unicode}{hyperref}
\PassOptionsToPackage{hyphens}{url}
\PassOptionsToPackage{dvipsnames,svgnames,x11names}{xcolor}
%
\documentclass[
  11pt,
  a4paper,
]{ltjsarticle}
\usepackage{amsmath,amssymb}
\usepackage{iftex}
\ifPDFTeX
  \usepackage[T1]{fontenc}
  \usepackage[utf8]{inputenc}
  \usepackage{textcomp} % provide euro and other symbols
\else % if luatex or xetex
  \usepackage{unicode-math} % this also loads fontspec
  \defaultfontfeatures{Scale=MatchLowercase}
  \defaultfontfeatures[\rmfamily]{Ligatures=TeX,Scale=1}
\fi
\usepackage{lmodern}
\ifPDFTeX\else
  % xetex/luatex font selection
\fi
% Use upquote if available, for straight quotes in verbatim environments
\IfFileExists{upquote.sty}{\usepackage{upquote}}{}
\IfFileExists{microtype.sty}{% use microtype if available
  \usepackage[]{microtype}
  \UseMicrotypeSet[protrusion]{basicmath} % disable protrusion for tt fonts
}{}
\makeatletter
\@ifundefined{KOMAClassName}{% if non-KOMA class
  \IfFileExists{parskip.sty}{%
    \usepackage{parskip}
  }{% else
    \setlength{\parindent}{0pt}
    \setlength{\parskip}{6pt plus 2pt minus 1pt}}
}{% if KOMA class
  \KOMAoptions{parskip=half}}
\makeatother
\usepackage{xcolor}
\usepackage[margin=24mm]{geometry}
\setlength{\emergencystretch}{3em} % prevent overfull lines
\providecommand{\tightlist}{%
  \setlength{\itemsep}{0pt}\setlength{\parskip}{0pt}}
\setcounter{secnumdepth}{5}
\ifLuaTeX
  \usepackage{selnolig}  % disable illegal ligatures
\fi
\IfFileExists{bookmark.sty}{\usepackage{bookmark}}{\usepackage{hyperref}}
\IfFileExists{xurl.sty}{\usepackage{xurl}}{} % add URL line breaks if available
\urlstyle{same}
\hypersetup{
  colorlinks=true,
  linkcolor={blue},
  filecolor={Maroon},
  citecolor={Blue},
  urlcolor={blue},
  pdfcreator={LaTeX via pandoc}}

\author{}
\date{}

\begin{document}

\hypertarget{chapter-04-ux5185ux7a4dux3068ux5e7eux4f55}{%
\section{Chapter 04 ---
内積と幾何}\label{chapter-04-ux5185ux7a4dux3068ux5e7eux4f55}}

ベクトル空間だけでは「近い」「垂直」「長い」という概念はありません。内積を導入して初めて幾何学が入ります。

\hypertarget{ux5185ux7a4dux304cux3082ux305fux3089ux3059ux3082ux306e}{%
\subsection{内積がもたらすもの}\label{ux5185ux7a4dux304cux3082ux305fux3089ux3059ux3082ux306e}}

\[
\langle x,y\rangle=x^Ty
\]

から

\begin{itemize}
\tightlist
\item
  ノルム \texttt{\textbar{}\textbar{}x\textbar{}\textbar{}}
\item
  距離 \texttt{\textbar{}\textbar{}x-y\textbar{}\textbar{}}
\item
  角度
\item
  直交
\item
  射影
\end{itemize}

が導かれます。

\hypertarget{ux306aux305cux76f4ux4ea4ux57faux5e95ux304cux7279ux5225ux304b}{%
\subsection{なぜ直交基底が特別か}\label{ux306aux305cux76f4ux4ea4ux57faux5e95ux304cux7279ux5225ux304b}}

一般基底で座標を求めるには連立方程式を解きます。正規直交基底なら

\[
c_i=q_i^Tx
\]

と内積一回で各成分を取り出せます。

つまり直交性は「幾何学的にきれい」なだけでなく、各方向の情報を独立に扱えるという計算上の意味があります。

\hypertarget{ux5c04ux5f71ux304bux3089ux6700ux5c0fux4e8cux4e57ux3078}{%
\subsection{射影から最小二乗へ}\label{ux5c04ux5f71ux304bux3089ux6700ux5c0fux4e8cux4e57ux3078}}

\texttt{b}
に最も近い部分空間上の点では、誤差ベクトルはその部分空間に垂直です。

この小学校的な「最短距離は垂線」という幾何学が

\[
A^T(b-Ax)=0
\]

という最小二乗の正規方程式になります。

\hypertarget{ux5230ux9054ux57faux6e96}{%
\subsection{到達基準}\label{ux5230ux9054ux57faux6e96}}

射影式を暗記するのでなく、「なぜ残差が直交すれば最短なのか」から導き直せることを目標にします。

\end{document}
